\chapter{Scientific Progress and Theory Choice}
\label{ch:progress}

\begin{chapterguide}
This chapter asks in what sense science progresses if theories are revised and
models are idealized. It distinguishes accumulation from improvement, considers
underdetermination and unconceived alternatives, and explains how \RCR\ treats
theory change as the enlargement and stabilization of constraint-sensitive
inference.
\end{chapterguide}

If objectivity is fallible, what makes one theory better than another? A later
theory may be less simple, less intuitive, or less familiar. It may reject
entities once considered real. It may also predict more accurately, explain
more phenomena, support new interventions, and clarify its own limitations.
Scientific progress is not the absence of change. It is change that increases
the reliability, range, precision, or depth of inquiry.

\section{Against simple accumulation}

The accumulation picture says that science adds facts to a fixed stock of
knowledge. Kuhn showed why this is inadequate: revolutions change concepts,
problems, and standards \citep{kuhn1962}. But the opposite picture, in which
each theory is a self-contained world, is also inadequate. Later science often
preserves measurements, limiting cases, mathematical relations, instruments,
and practical controls.

\RCR\ identifies five forms of progress:

\begin{enumerate}
\item \textbf{Empirical progress:} more reliable phenomena are established and
fewer important results remain unexplained.
\item \textbf{Predictive progress:} the theory yields new, risky, and calibrated
predictions.
\item \textbf{Causal progress:} the theory identifies more stable mechanisms
and intervention boundaries.
\item \textbf{Representational progress:} models preserve relevant structure
across wider ranges or provide better approximations.
\item \textbf{Reflexive progress:} the inquiry becomes better at identifying
its own errors, assumptions, and social exclusions.
\end{enumerate}

These dimensions can conflict. A more accurate model may be less interpretable.
A broader theory may weaken local fit. A new measure may reveal that an old
effect was an artifact. Progress is therefore multidimensional and must be
reported as a profile.

\section{The problem of underdetermination}

For any finite set of observations, many models may fit. This is
underdetermination. It follows from the fact that data constrain hypotheses
but do not usually select one uniquely. Duhem's holism and Quine's web of
belief make the problem vivid \citep{duhem1906,quine1951}.

The existence of alternatives does not imply that theory choice is arbitrary.
Researchers can generate new constraints by measuring a different variable,
designing a discriminating experiment, comparing interventions, or testing
novel domains. The appropriate response to underdetermination is not to
declare victory for the preferred theory. It is to ask what further test would
separate the live alternatives and whether the difference matters to the
question.

Kyle Stanford's discussion of unconceived alternatives strengthens the
warning. Even when scientists compare the theories they have imagined, they
may fail to consider a future theory that preserves present successes while
reinterpreting their ontology \citep{stanford2006}. \RCR\ accepts that no
current commitment is immune to future alternatives. It responds through
graded commitment and a distinction between secure constraints and provisional
interpretations.

\section{Convergence and robustness}

Convergence is not mere agreement. Two studies that use the same flawed
instrument and code may be correlated failures. Convergence is epistemically
strong when the routes share the target but differ in potential errors.

Examples of productive convergence include:

\begin{itemize}
\item a physical quantity measured by distinct instrument principles;
\item a biological mechanism supported by molecular, cellular, and organismic
interventions;
\item a climate attribution supported by observations, physical models,
energy-balance reasoning, and independent indicators;
\item a policy effect supported by randomization, quasi-experiment, mechanism,
and implementation evidence.
\end{itemize}

Robustness tests should be chosen to vary what is not supposed to matter and
to preserve what is supposed to matter. If a claim survives only when every
analytic choice is held fixed, its objectivity is weaker than its apparent
precision suggests.

\section{Simplicity and complexity}

Simplicity is often treated as a reason to prefer one theory. But simple in
what sense? A theory can have fewer parameters but a more complicated
ontology. A concise model can hide a complex data-processing pipeline. A
complex model can be more honest if the target is heterogeneous.

\RCR\ treats simplicity as a resource for transport and criticism, not as a
metaphysical law. A model is preferable when its complexity is proportionate
to the structure it must represent and when the added complexity earns new
constraint-sensitive success. This is close to the methodological idea behind
penalizing overfit models, but it is broader: complexity includes hidden
assumptions, institutional dependencies, and unreported flexibility.

\section{The role of explanation}

Scientific progress is not only increased prediction. A black-box system can
forecast accurately but fail to reveal which intervention would change the
outcome. An explanation makes a difference to what questions can be asked next.

Philip Kitcher connects scientific progress with the organization and
unification of knowledge \citep{kitcher1993}. A good explanation can reduce
dependence on isolated facts and improve the ability to derive new consequences.
But unification can be forced. \RCR\ counts explanatory reach only when the
bridges between domains are explicit and survive independent tests.

\section{Theory change and retained constraints}

A theory change is progressive when it replaces a representation while
retaining or improving constraints that mattered. Consider a sequence of
models \(M_1,M_2,\ldots,M_k\). A progress relation is not simply
\(M_{j+1}\) has more equations than \(M_j\). It can include:

\[
  \operatorname{Prog}(M_{j+1},M_j)
  =
  \operatorname{Retain}(R)
  +
  \operatorname{Extend}(D)
  +
  \operatorname{Correct}(A),
\]
where \(R\) is a set of retained relations, \(D\) is newly constrained domain
behaviour, and \(A\) is a set of known anomalies or assumptions corrected by the
new model. The expression is conceptual, not a numerical theorem.

A later theory can therefore be better even if it says that an older entity was
not fundamental. It may retain the old theory as a limiting approximation,
explain why the approximation worked, and predict where it fails. This is one
meaning of objective progress.

\section{Research programmes and the future}

A scientific field should not aim only to confirm its current framework. It
should cultivate a portfolio of activities:

\begin{enumerate}
\item normal work that improves precision and solves established puzzles;
\item anomaly investigation that tests the framework's boundaries;
\item alternative model construction that exposes hidden assumptions;
\item intervention and replication that probe causal stability;
\item conceptual analysis that revises categories when measurements conflict;
\item institutional reform that improves correction capacity.
\end{enumerate}

A field that lacks anomaly work can become dogmatic. A field that lacks normal
work cannot stabilize knowledge. A field that lacks alternatives cannot know
whether its theory is overfitted. A field that lacks institutional reform can
identify errors without being able to correct them.

\begin{principlebox}
\textbf{Progress principle.} Scientific progress is the cumulative enlargement
and stabilization of world-sensitive constraints, accompanied by improved
ability to identify scope, error, and alternative representations. It need not
preserve a fixed ontology or produce a final theory.
\end{principlebox}

\section*{Chapter summary}

Progress is multidimensional: empirical, predictive, causal,
representational, and reflexive. Underdetermination is real but can be reduced
by discriminating tests and independent routes of evidence. Convergence requires
different potential error structures, not repeated agreement from one pipeline.
Simplicity is a methodological resource rather than a metaphysical guarantee.
Theory change is progressive when it retains important constraints, extends
their domain, corrects anomalies, and improves the ability to find errors.

\begin{discussion}
\begin{enumerate}
\item Can a theory be progressive if it increases prediction but decreases
interpretability?
\item How can researchers search for alternatives they have not yet imagined?
\item Which forms of progress matter most in policy-relevant science?
\end{enumerate}
\end{discussion}

\begin{furtherreading}
See Lakatos (1978), Laudan (1977), Kitcher (1993), Stanford (2006),
Cartwright (1999), and the literature on structural realism and scientific
progress.
\end{furtherreading}

